\subsection{Ступени полёта}

Каждая программа ниже --- на одну ступень сложнее предыдущей. Запускайте каждую, смотрите камерой \ui{Дрон} и читайте журнал на вкладке \ui{Вывод}.

\subsubsection*{Ступень 1. Только моторы}

\begin{codefile}{lua/06_missions/missions_engines_arm.lua}
\end{codefile}

Одна команда --- подготовка: автопилот запускает моторы, дрон остаётся на земле, а подтверждение приходит событием \verb|ENGINES_STARTED|. Так проверяют моторы на настоящем дроне со снятыми пропеллерами. В справочнике есть и команда \verb|ap.push(Ev.ENGINES_ARM)| (приложение~\ref{sec:lua-api}); в симуляторе она делает то же, что \verb|MCE_PREFLIGHT|, и повторять её после подготовки не нужно --- автопилот уже на другом этапе и отклонит её. Без взлёта подготовка через 8~с истекает, и моторы останавливаются.

\subsubsection*{Ступень 2. Взлёт и посадка}

\begin{codefile}{lua/06_missions/mission_takeoff_landing.lua}
\end{codefile}

\paragraph{Что происходит.} Предстарт $\to$ через секунду взлёт $\to$ по событию \verb|TAKEOFF_COMPLETE| посадка $\to$ по \verb|COPTER_LANDED| команда \verb|ENGINES_DISARM|. После посадки автопилот и сам выключает моторы, так что эта строка --- страховка; зато она показывает единственное безопасное место для \verb|ENGINES_DISARM| в штатной миссии. Команда посадки только \emph{начинает} снижение; закончилась посадка, когда пришло \verb|COPTER_LANDED|. Четыре строки в \verb|callback| --- это вся миссия.

\begin{center}
\simshot[0.85\linewidth]{05-flight.png}
\captionof{figure}{Камера \ui{Дрон} следует за дроном; голубая линия --- след полёта.}
\end{center}

\subsubsection*{Ступень 3. Полёт по прямой}

\begin{codefile}{lua/06_missions/missions_straight_flight.lua}
\end{codefile}

После взлёта --- точка $(1.0, 0.0, 0.8)$: метр по оси X (поперёк носа) на высоте 0.8~м. Точку отправляют только после \verb|TAKEOFF_COMPLETE|: команда движения во время взлёта --- ошибка «Маршрут слишком рано». Посадка --- по \verb|POINT_REACHED|, то есть когда дрон действительно там.

\begin{warnbox}[Посадка раньше точки]
Частая ошибка --- отправить точку и в том же обработчике, через таймер на пару секунд, --- посадку. По умолчанию дрон летит со скоростью около 0.4~м/с (параметр \verb|Copter_pos_vMax|): за две секунды он не долетит и метра, и посадка отменит перелёт. Следующая команда --- по следующему событию.
\end{warnbox}

\begin{taskbox}[Одна точка: оси и знаки]
Меняйте в \texttt{missions\_straight\_flight.lua} только цель и каждый раз предсказывайте результат до запуска.
\begin{enumerate}
  \item Цель $(0, 0, 2)$: меняется только Z.
  \item Цель $(1, 1, 1)$. Затем $(-1, 1, 1)$: изменился только знак X.
  \item Цель $(2, 0, 1)$; перед посадкой напечатайте \verb|"Прибыл"| и положение из \verb|Sensors.lpsPosition()|.
\end{enumerate}
\expected{1) дрон поднимается над стартом до 2~м и садится, не смещаясь по земле; 2) вторая точка зеркальна первой относительно оси Y --- дрон уходит в другую сторону по X; 3) «Прибыл» появляется только после остановки у цели, а напечатанный X близок к 2, но не равен ему точно (например, 1.87 или 2.04).}
\mission{Lua}{38, 41, 44}
\end{taskbox}

\subsubsection*{Ступень 4. Несколько точек и счётчик шагов}

Каждая точка заканчивается одним и тем же событием \verb|POINT_REACHED|. Чтобы знать, \emph{какая} точка достигнута, нужен счётчик:

\begin{codefile}{lua/06_missions/missions_points_flight.lua}
\end{codefile}

\paragraph{Что происходит.} Маршрут --- таблица точек, \verb|curr| --- номер следующей. Функция \verb|nextPoint| отправляет точку с номером \verb|curr| и увеличивает счётчик; когда точки кончились --- садится. Её вызывают два события: \verb|TAKEOFF_COMPLETE| (первая точка) и \verb|POINT_REACHED| (все остальные). Координаты отделены от логики полёта: чтобы изменить маршрут, правят только таблицу. Это основной шаблон маршрута в курсе; так же устроены официальные примеры Геоскан. Отправить вторую точку сразу вслед за первой нельзя: новая цель заменит прежнюю, и первую дрон не посетит.

\subsubsection*{Ступень 5. Фигуры из точек}

Треугольник --- три вершины на окружности радиуса $r$ под углами $0^\circ$, $120^\circ$, $240^\circ$: $x = r\cos a$, $y = r\sin a$.

\begin{codefile}{lua/06_missions/missions_triangle_flight.lua}
\end{codefile}

Квадрат со стороной $a$ и центром в начале координат --- вершины $(\pm a/2, \pm a/2, z)$:

\begin{codefile}{lua/06_missions/missions_square_flight.lua}
\end{codefile}

\begin{hintbox}[Время перелёта]
У \verb|ap.goToLocalPoint| есть четвёртый аргумент --- время перелёта в секундах: \verb|ap.goToLocalPoint(1, 0, 0.8, 6)| долетит за 6~с. Без него скорость берётся из параметров автопилота.
\end{hintbox}

\begin{taskbox}[Маршрут по этапам]
\begin{enumerate}
  \item Угол: через вершину $(0, 2, 1)$ в $(2, 2, 1)$. Вершину нужно посетить, а не срезать.
  \item Квадрат $(2, 0, 1) \to (2, 2, 1) \to (0, 2, 1) \to (0, 0, 1)$ таблицей точек в \texttt{missions\_points\_flight.lua}.
  \item Возвращение домой: до подготовки запомните \verb|local homeX, homeY = Sensors.lpsPosition()|, посетите $(0, 2, 1)$ и вернитесь в \verb|(homeX, homeY, 1)|, затем сядьте. Проверьте, перенеся дрон в другое место сцены.
  \item Свой маршрут из четырёх точек задания «Разведка площадки» с возвратом домой. Запустите его три раза подряд, не меняя код.
\end{enumerate}
\expected{1) след на сцене --- прямой угол с вершиной в $(0, 2)$; 2) замкнутый квадрат, посадка в начале сцены; 3) дрон садится там, откуда взлетел, даже после переноса, а не в начале сцены; 4) три одинаковых следа: каждый запуск начинается со сброса сцены, поэтому из одинаковых условий маршрут повторяется.}
\mission{Lua}{46--54}
\end{taskbox}

\subsection{Полигоны и местность}

Пустая площадка прощает всё. На полигоне маршрут проверяется местностью: деревьями, холмами, зданиями, мачтами, воротами. Готовые полигоны --- в панели \ui{Сцена}, раздел «Готовые полигоны»: учебные лес, холмы, город, промышленная зона, экспедиция (смешанная местность), гоночная трасса и другие. В миссиях 55--63 нужная учебная сцена создаётся сама; её подпись --- «Учебная сцена · урок~$N$ · seed~$S$».

\begin{itemize}
    \item \textbf{Seed} --- число, из которого генератор получает расстановку объектов. Тот же seed --- та же карта, поэтому трудный маршрут можно повторять и отлаживать. Другой seed --- другая расстановка, и старые координаты обхода могут стать неверными: сначала осмотрите карту.
    \item \textbf{Размер дрона.} Координаты цели --- это центр дрона, а винты выступают в стороны. Коридор между деревьями должен быть шире корпуса с запасом, стартовая зона --- свободна вокруг винтов.
    \item \textbf{Высота до препятствия.} Высоту поднимают \emph{до} того, как маршрут пересечёт холм или мачту: отдельной точкой над стартом, а не на ходу. Высоту препятствия узнают \key{Ctrl}+щелчком по нему.
    \item \textbf{Обход.} Если прямая к цели пересекает здание, путь строят из нескольких точек с запасом от фасада. Ворота проходят тремя точками --- подход, центр проёма, выход, --- иначе дрон зацепит раму, срезая угол.
    \item \textbf{Участки.} Длинный маршрут через разные зоны делят на участки, у каждого --- своя высота и свои ограничения; так его легче проверять по частям.
\end{itemize}

\begin{taskbox}[Маршрут по местности]
Откройте полигон «Учебный: лес» и проложите маршрут на 8~м вперёд по Y через свободный коридор: подъём над стартом, точка в конце коридора, посадка. Затем на полигоне «Учебный: холмы» пройдите 12~м вперёд по Y, подняв высоту до начала пересечения холмов. Высоты выбирайте по \key{Ctrl}+щелчку.
\expected{в лесу след идёт по коридору, не задевая крон, событие \texttt{SHOCK} не приходит; на холмах дрон набирает высоту над стартом, потом летит горизонтально, и дальномер \texttt{Sensors.range()} в самой высокой точке холма показывает просвет больше нуля.}
\mission{Lua}{55--63}
\end{taskbox}

\subsection{Измерения: цель, положение, допуск}

Команда сообщает, куда дрон \emph{должен} прилететь; \textbf{телеметрия} --- где он \emph{на самом деле}. Это разные числа, и проверять миссию нужно по второму.

\begin{codefile}{lua/06_missions/missions_telemetry.lua}
\end{codefile}

\paragraph{Что происходит.}
\begin{itemize}
    \item \verb|report| печатает \emph{измеренное} положение с подписью: до подготовки, после взлёта, по прибытии, после посадки. Цель печатается отдельной строкой --- в журнале видно и намерение, и результат. Такой журнал --- отчёт полёта: по нему находят первый этап, где поведение отклонилось.
    \item Таймер раз в 0.5~с печатает \textbf{ошибку} --- цель минус измерение. Знак ошибки показывает, в какую сторону ещё лететь; к концу перелёта она уменьшается к нулю. Период опроса выбирают по задаче: слишком редкий пропустит быстрые изменения, слишком частый засорит журнал.
    \item Прибытие проверяют с \textbf{допуском}: \verb|math.abs(target[1] - x) <= tolerance|. Измерение и управление неточны, поэтому \verb|x == 2| почти никогда не выполняется.
    \item \textbf{Тайм-аут} ограничивает ожидание и называет причину остановки: «тайм-аут: цель не достигнута\dots» --- не то же самое, что штатное «Цель достигнута».
\end{itemize}

Ещё три приёма из миссий 66--70:
\begin{itemize}
    \item \textbf{Нет данных.} Отсутствующее измерение --- \verb|nil|, а не ноль. Сначала \verb|if value == nil then print("Нет данных") return end|, и только потом сравнение с числом, иначе --- ошибка «Сравнение с пустым значением». Координаты в симуляторе есть всегда, но так пишут любую обработку датчика.
    \item \textbf{Зона спокойствия.} Возле границы шум измерений переключает решение туда и обратно. Ошибку внутри \verb|[-tolerance, tolerance]| считают нулём: решение «лететь дальше» принимают, только когда ошибка явно больше допуска.
    \item \textbf{Ограничение.} Рассчитанная величина --- скорость, шаг, яркость --- должна оставаться в допустимом диапазоне: \verb|math.max(-limit, math.min(limit, v))|. Большая ошибка не означает неограниченную скорость.
\end{itemize}

\begin{taskbox}[Отчёт по измерениям]
\begin{enumerate}
  \item Запустите пример и выпишите из журнала строки «Цель» и «Прибытие». Совпадают ли числа?
  \item Замените период опроса 0.5~с на 0.2~с, потом на 2~с. Сколько строк «Ошибка X» получается за перелёт?
  \item Уменьшите \verb|timeout| до 3~с. Какая причина остановки появится в журнале?
  \item Напишите функцию \verb|zone(e, tol)|: 0 --- если \verb|e| внутри допуска, 1 --- больше допуска, $-1$ --- меньше минус допуска. Проверьте на ошибках $-0.04$, $0.03$, $-0.02$, $0.2$, $-0.2$ при допуске 0.1.
\end{enumerate}
\expected{1) не совпадают точно: прибытие отличается от цели на сантиметры, но в пределах 0.15~м, поэтому печатается «Цель достигнута»; 2) при 0.2~с строк примерно в 10 раз больше, чем при 2~с, а ошибка в них убывает к нулю; 3) «тайм-аут: цель не достигнута за 3 с», и дрон садится, не долетев; 4) \texttt{0 0 0 1 -1}.}
\mission{Lua}{64--72}
\end{taskbox}
